\subsection{Variétés allotropiques du soufre}

Le soufre existe à l’état solide sous deux variétés $\ce{S_8 (\alpha)}$ (orthorhombique) et $\ce{S_8 (\beta)}$ (monoclinique). Les volumes molaires sont $V^\star_{m,\alpha} = 15.3 ~ \si{\centi\meter\cubed\per\mole}$ pour le soufre $\ce{S_8 (\alpha)}$ et $V^\star_{m,\alpha} = 16.4 ~ \si{\centi\meter\cubed\per\mole}$ pour $\ce{S_8 (\beta)}$. Sous une pression $P=P^\circ$, le point de transition réversible entre les deux phases est $\theta_{tr} = 95.5 ~ \si{\celsius}$. On suppose les solides incompressibles. Déterminer l’accroissement $\Delta T$ de la température de transition réversible consécutif à une augmentation de pression de \SI{10.0}{\bar}.

\textbf{Donnée :} $S_m^\star (\beta) - S_m^\star (\alpha) \approx  S_m^{\circ, \star} (\beta) - S_m^{\circ, \star} (\alpha) = 7.87 ~ \si{\joule\per\mole\per\kelvin}$